%=============================================================================!
% 18.5  GRADIENT DESCENT                                                      !
%=============================================================================!
\section{Gradient descent}
\label{sec:lo-descent}

A walker caught in fog on a hillside, wanting the valley floor, feels the
slope under their feet and steps downhill, then feels again. Short steps
are safe but slow; a long step on a steep slope can carry the walker
across a narrow valley and up the far side, higher than they started,
and a run of such steps climbs out of the valley altogether. The normal
equations solved the fits so far in one calculation, but most models,
the networks of planned Chapter~19 among them, have losses that no formula
minimises, and they are walked down instead.

\subsection*{The step and its limit}

\emph{Gradient descent} repeats the step
\begin{equation}
   \vb{w} \leftarrow \vb{w} - \eta\,\grad L(\vb{w}) ,
   \label{eq:lo-descent}
\end{equation}
moving the weights against the gradient, the direction in which the loss
falls fastest (\cref{sec:en-gradient}), by an amount set by the
\emph{learning rate} $\eta$. Near a minimum with positive curvature in
every direction, a smooth loss is approximated to second order by a
quadratic bowl, just as a potential energy is near a stable minimum
(\cref{sec:os-modes}). In this approximation, with $\vb{e} = \vb{w} - \vb{w}^\ast$ the
departure from the minimum $\vb{w}^\ast$ and $A$ the matrix of second
derivatives there, the Hessian,
\begin{equation*}
   L = L^\ast + \tfrac12\vb{e}^{\mathsf T}A\vb{e} , \qquad
   \grad L = A\vb{e} .
\end{equation*}
For this quadratic, a step changes the departure by
\begin{align*}
   \vb{e} &\leftarrow \vb{e} - \eta A\vb{e} = (I - \eta A)\,\vb{e}
      && \text{subtracting $\vb{w}^\ast$ from both sides of the step}\\
   c_k &\leftarrow (1 - \eta\lambda_k)\,c_k
      && \text{along each eigenvector of $A$, of eigenvalue $\lambda_k$,}
\end{align*}
the components $c_k$ of $\vb{e}$ along the eigenvectors changing
independently, as the normal modes of a vibrating molecule do. Each
shrinks only if $\lvert1 - \eta\lambda_k\rvert < 1$, that is $0 < \eta <
2/\lambda_k$, and the steepest direction, with the largest eigenvalue
$\lambda_{\max}$, sets the limit:
\begin{equation}
   \eta < \frac{2}{\lambda_{\max}} .
   \label{eq:lo-limit}
\end{equation}
This is the stability limit of an integrator in another guise
(\cref{sec:in-stability}): the fastest vibration of a molecule bounds the
time step, and the steepest direction of a loss bounds the learning rate.
Between $1/\lambda_{\max}$ and $2/\lambda_{\max}$ the factor $1 -
\eta\lambda_{\max}$ is negative, and the steep component flips sign at
every step while it shrinks, the zig-zag of the walker across the valley.

\Cref{fig:lo-descent} walks a loss of two weights whose Hessian has the
eigenvalues 1 and 10, so $2/\lambda_{\max} = 0.2$. At $\eta = 0.05$ the
factors are 0.95 and 0.50, and the shallow direction crawls; at 0.18 they
are 0.82 and $-0.80$, and the path zig-zags down the valley; at 0.205
the steep factor is $-1.05$, and the path climbs out. The best fixed rate
makes the slowest factor as small as it can be, which it is when the
two extreme factors are equal and opposite, $1 - \eta\lambda_{\min} =
-(1 - \eta\lambda_{\max})$:
\begin{equation*}
   \eta = \frac{2}{\lambda_{\min} + \lambda_{\max}} , \qquad
   1 - \eta\lambda_{\min} = \frac{\lambda_{\max} - \lambda_{\min}}
   {\lambda_{\max} + \lambda_{\min}} = \frac{\kappa - 1}{\kappa + 1} ,
\end{equation*}
with $\kappa = \lambda_{\max}/\lambda_{\min}$ the \emph{condition number}
of the Hessian. Here $\eta = 0.182$, the factor is 0.818, and since the
loss falls as the square of the factor, it falls a millionfold in
$\ln10^6/(-2\ln0.818) = 34$ steps.

\subsection*{The condition number decides}

For a large condition number the factor is $1 - 2/(\kappa + 1)$, and the
steps for a millionfold fall grow in proportion to $\kappa$, about
$\kappa\ln10^6/4$. The least-squares loss has the Hessian $2X^{\mathsf
T}X$, and the polynomial fits of \cref{sec:lo-basis} are badly
conditioned: $X^{\mathsf T}X$ has $\kappa = 74$ for degree 3 in the scaled
distance, \num{4.6e6} for degree 8 and \num{5.1e16} for degree 14, and
for degree 8 in $r$ itself \num{3.3e26}, beyond what the 16 digits of
double precision can even resolve. Gradient descent at its best fixed rate
would need about 260 steps for degree 3 and \num{1.6e7} for degree 8.
Scaling the variable, as \cref{sec:lo-basis} did, is the first remedy;
the solvers of the normal equations are another; and the optimisers of
\cref{sec:lo-momentum} are a third, for the losses that have no normal
equations.

A loss is \emph{convex} if the straight line between any two points of
its graph lies on or above the graph; a quadratic with a Hessian whose
eigenvalues are all positive or zero is convex, and the least-squares
loss, with Hessian $2X^{\mathsf T}X$, is one. A convex loss has no local
minima besides its least value, so gradient descent at a stable rate
reaches it from any start. The losses of neural networks are not convex,
and where gradient descent ends then depends on where it began
(planned Chapter~19).

\begin{figure}[htb]
\centering
\includegraphics{ch18_learning/descent}
\caption{Gradient descent on a quadratic loss of two weights whose
Hessian has the eigenvalues 1 and 10 along axes turned by \ang{30}. (a)
The paths at learning rates of 0.05, 0.18 and 0.205 over the contours of
the loss (first 25 steps, 9 for the rate that diverges). (b) The loss
against the step for rates of 0.05, 0.1, 0.18 and 0.205, with the term of
each rate's slowest direction, $\frac12\lambda c^2(1 - \eta\lambda)^{2t}$
for $c$ the start's component along it (dashed).}
\label{fig:lo-descent}
\end{figure}

\begin{keyidea}
Gradient descent steps against the gradient by the learning rate $\eta$.
Near the minimum each eigen-direction of the Hessian shrinks by $1 -
\eta\lambda$ a step, so $\eta < 2/\lambda_{\max}$, the stability limit of
an integrator again; at the best fixed rate the slowest direction
shrinks by $(\kappa - 1)/(\kappa + 1)$, and the condition number $\kappa$
sets how many steps are needed. A convex loss has one least value, which
descent at a stable rate reaches.
\end{keyidea}

\notebook{18}{walk down a valley at different rates and condition numbers}

\begin{selfcheck}
A loss has a Hessian with eigenvalues 0.5 and 50. What is the largest
stable learning rate, and what factor does the slow direction shrink by
at the best fixed rate?
\end{selfcheck}
